% Part: second-order-logic
% Chapter: syntax-and-semantics
% Section: terms-formulas

\documentclass[../../../include/open-logic-section]{subfiles}

\begin{document}

\olfileid{sol}{syn}{frm}

\olsection{الحدود و\printtoken{P}{formula}}

يبدأ إنشاء تعبيرات منطق الرتبة الثانية، كما في الأولى، بمعجم
أصلي من \emph{!!{variable}s} و\emph{!!{constant}s}
و\emph{!!{predicate}s}، وقد يضم \emph{!!{function}s}.
ثم تؤلف منه، مع الروابط والمكممات والأقواس والفواصل وغيرها من
علامات الترقيم، \emph{الحدود} و\emph{!!{formula}s}.
والزيادة في الرتبة الثانية أن المتغيرات لا تقتصر على الأفراد؛
بل تتناول العلاقات والدوال أيضًا، ويجوز التكميم عليها.
فالرموز المنطقية هي رموز الرتبة الأولى مع ما يأتي:

\begin{enumerate}
\item لكل عدد مواضع $\OLId{latin-ordinary}{n}$ مجموعة !!{denumerable} من
  !!{variable}s العلاقات من الرتبة الثانية:
  $\Obj V_\OLMathNumeral{0}^\OLId{latin-ordinary}{n}$، و$\Obj V_\OLMathNumeral{1}^\OLId{latin-ordinary}{n}$، و$\Obj V_\OLMathNumeral{2}^\OLId{latin-ordinary}{n}$، و\dots
\item وللدوال من الرتبة الثانية مجموعة !!{denumerable} من
  !!{variable}s: $\Obj u_\OLMathNumeral{0}^\OLId{latin-ordinary}{n}$، و$\Obj u_\OLMathNumeral{1}^\OLId{latin-ordinary}{n}$، و$\Obj u_\OLMathNumeral{2}^\OLId{latin-ordinary}{n}$، و\dots
\end{enumerate}

اصطلحنا على $\OLId{latin-ordinary}{x}$ و$\OLId{latin-ordinary}{y}$ و$\OLId{latin-ordinary}{z}$ متغيرات فوقية لمتغيرات الرتبة
الأولى $\Obj v_\OLId{latin-ordinary}{i}$. وعلى المثال نفسه نستعمل $\OLId{latin-looped}{X}$ و$\OLId{latin-looped}{Y}$ و$\OLId{latin-looped}{Z}$
وما شابهها متغيرات فوقية لـ$\Obj V_\OLId{latin-ordinary}{i}^\OLId{latin-ordinary}{n}$، ونستعمل $\OLId{latin-ordinary}{u}$ و$\OLId{latin-ordinary}{v}$
وما شابههما لـ$\Obj u_\OLId{latin-ordinary}{i}^\OLId{latin-ordinary}{n}$.

\begin{explain}
وأما الرموز غير المنطقية للغة الرتبة الثانية فتتحدد كما في
الأولى: نعيّن قائمة !!{constant}s و!!{function}s
و!!{predicate}s الخاصة بها.

وتضم لغة الرتبة الأولى عادةً الـ!!{identity} $\eq$.
وفي الرتبة الأولى لا غنى عن الرموز غير المنطقية للغة $\Lang{L}$
للتعبير عن أمور ذات شأن؛ فنعم، توجد !!{sentence}s خالية
منها، لكن الاقتصار على $\eq$ يضيّق ما يمكن قوله.
وأما الرتبة الثانية فتتيح عددًا غير محدود من متغيرات العلاقات
والدوال؛ فبها يمكن التعبير عما نعبر عنه في لغة الرتبة الأولى
من غير أن نخص اللغة بمجموعة من الرموز غير المنطقية.
\end{explain}

\begin{defn}[حدود الرتبة الثانية]
مجموعة \emph{حدود الرتبة الثانية} في $\Lang L$، ورمزها
$\TrmSOL[\OLId{latin-looped}{L}]$، هي المجموعة التي يعرّفها
\olref[fol][syn][frm]{defn:terms} بعد إلحاق الشرط الآتي به:
\begin{enumerate}
\item متى كان $\OLId{latin-ordinary}{u}$ متغيرًا داليًا ذا $\OLId{latin-ordinary}{n}$ مواضع،
وكانت $\OLId{latin-ordinary}{t}_\OLMathNumeral{1}$، \dots، $\OLId{latin-ordinary}{t}_\OLId{latin-ordinary}{n}$ حدودًا، عُدّ
$\Atom{u}{t_1, \ldots, t_n}$ حدًا.
\end{enumerate}
\end{defn}

\begin{explain}
فصورة حد الرتبة الثانية صورة حد الرتبة الأولى، سوى أن الموضع
الذي تشغله !!a{function} $\Obj{f^n_i}$ في الأول قد يشغله
المتغير الدالي $\Obj{u^n_i}$ في الثاني.
\end{explain}

\begin{defn}[\usetoken{s}{formula} الرتبة الثانية]
نحصل على مجموعة \emph{!!{formula}s الرتبة الثانية}
$\FrmSOL[\OLId{latin-looped}{L}]$ للغة $\Lang L$ بإلحاق الشروط الآتية
بـ\olref[fol][syn][frm]{defn:formulas}:
\begin{enumerate}
\item متى كان $\OLId{latin-looped}{X}$ متغيرًا محموليًا ذا $\OLId{latin-ordinary}{n}$ مواضع،
وكانت $\OLId{latin-ordinary}{t}_\OLMathNumeral{1}$، \dots، $\OLId{latin-ordinary}{t}_\OLId{latin-ordinary}{n}$ حدودًا من الرتبة الثانية في
$\Lang L$، عُدّت $\Atom{X}{t_1,\ldots, t_n}$ !!{formula} ذرية.

\tagitem{prvAll}{من !!a{formula} $!A$ ومتغير دالي $\OLId{latin-ordinary}{u}$
  تتكون !!a{formula} $\lforall[\OLId{latin-ordinary}{u}][!A]$.}{}

\tagitem{prvAll}{من !!a{formula} $!A$ ومتغير محمولي $\OLId{latin-looped}{X}$
  تتكون !!a{formula} $\lforall[\OLId{latin-looped}{X}][!A]$.}{}

\tagitem{prvEx}{من !!a{formula} $!A$ ومتغير دالي $\OLId{latin-ordinary}{u}$
  تتكون !!a{formula} $\lexists[\OLId{latin-ordinary}{u}][!A]$.}{}

\tagitem{prvEx}{من !!a{formula} $!A$ ومتغير محمولي $\OLId{latin-looped}{X}$
  تتكون !!a{formula} $\lexists[\OLId{latin-looped}{X}][!A]$.}{}
\end{enumerate}
\end{defn}

\end{document}
